%! Random Sampling Methods — Unit 3, Lesson 3.2 — Table Activity: Which Plan Fits? — KEY
% ONE-OFF, STANDALONE HANDOUT. Not a packet component: `table_activity` is deliberately NOT a
% name in shared/lesson.mk's STUDENT_ORDER, so `make` never builds it and no packet ever
% merges it. The teacher prints it only on a day it is used (see the optional box at the end
% of ../main.tex). Build it by hand, out of tree:
%   cd unit03/lesson02/table_activity_key && TEXINPUTS="$(git rev-parse --show-toplevel)/shared//:" \
%     latexmk -xelatex -outdir=../../../target/unit03/lesson02/table_activity_key main.tex
%
% THREE TABLES, THREE CARDS, ONE BEST METHOD EACH — each sheet is exactly two pages
% (print duplex, one sheet per student at that table):
%   Table 1  Pine Ridge County smoke detectors  -> CLUSTER   (no list, in-person visits,
%            map sections each a mix like the county)
%   Table 2  StreamLine music app ratings       -> SRS       (full numbered list, free to reach
%            anyone, nothing known to stratify on)
%   Table 3  Northfield Foods commute distance  -> STRATIFIED by work site (commutes alike
%            inside a site) — and THE CRUX: "pick one site as a cluster"
% Contexts are new: not Lakeview / Riverside (notes), the hospital (homework), or farms /
% textbooks / the airplane (AP Practice).
%
% PAGE LOCKSTEP with ../table_activity/: the blank is this file with every \answerspace,
% \cellans and \fillin second argument emptied — nothing else differs but the style package
% and the header line.
\documentclass[12pt]{article}
\usepackage{apstats-article}
\usepackage{apstats-key}
\newcommand{\answerspace}[2]{\par\nopagebreak\noindent\begin{minipage}[t][#1][t]{\linewidth}%
  \color{keyred}\bfseries\footnotesize #2\end{minipage}\par}
\newcommand{\cellans}[2]{\begin{minipage}[t][#1][t]{\linewidth}\raggedright
  \color{keyred}\bfseries\footnotesize #2\end{minipage}}
\newcommand{\fillin}[2]{\underline{\makebox[#1][l]{\ \smash{\color{keyred}\bfseries #2}}}}
\newcolumntype{L}[1]{>{\raggedright\arraybackslash}p{#1}}

\begin{document}

% ════════════════════════════════════════════════════════════════════════════
% TABLE 1 — PINE RIDGE COUNTY — CLUSTER
% ════════════════════════════════════════════════════════════════════════════
\pageheader{Unit 3, Lesson 3.2}{Table Activity: Which Plan Fits? \quad --- \quad Table 1 \quad --- \quad Key}

\boxguard[10]
\begin{headlinebox}{sky}
\small
\textbf{Your table is the survey team.} Read your card, try all three sampling methods on it,
and choose the one that fits. Be ready to defend your choice to the class in one minute ---
\emph{in context}, using the words \emph{alike inside}, \emph{looks like the whole}, or
\emph{cost}.
\end{headlinebox}

\vspace{0.08in}

\boxguard[12]
\begin{scenariobox}[Table 1 card: Smoke Detectors in Pine Ridge County]{navy}
\small
\textbf{Statistical question:} What proportion of homes in Pine Ridge County have a working
smoke detector?
\begin{itemize}[leftmargin=1.3em, itemsep=2pt, topsep=4pt]
  \item About 9{,}000 homes are spread over 600 square miles of farmland, hills, and small
        towns.
  \item \textbf{There is no list of homes.} The county map does divide the county into 150
        \emph{map sections} of about 60 homes each.
  \item Every map section is a mix: old and new houses, owners and renters, families of every
        income.
  \item A fire inspector must visit each home \emph{in person} and press the detector's test
        button. Driving from one map section to another can take an hour.
  \item The fire department can afford to check about \textbf{300 homes}.
\end{itemize}
\end{scenariobox}

\vspace{0.08in}

\boxguard[10]
\begin{scenariobox}[1. Read the population]{navy}
\small
\begin{enumerate}[label=\textbf{\alph*.}, leftmargin=1.6em, itemsep=4pt]
  \item What are the individuals? What is the variable?
        \answerspace{1.1cm}{Individuals: the homes in Pine Ridge County. Variable: whether the
        home has a working smoke detector (categorical: yes or no).}
  \item Is there a list of every individual? What does it cost to reach one?
        \answerspace{1.3cm}{No --- there is no list of homes. Every home needs an in-person
        visit, and the driving between map sections is the big cost.}
  \item Are there groups that are \emph{alike inside} on the response? Groups that each
        \emph{look like the whole} population?
        \answerspace{1.3cm}{The map sections each look like the whole county --- every one is a
        mix of old and new, owners and renters. Nothing tells us which homes are alike inside.}
\end{enumerate}
\end{scenariobox}

\clearpage

\boxguard[16]
\begin{scenariobox}[2. Try all three methods on your card]{navy}
\small
\renewcommand{\arraystretch}{1.2}
\begin{tabular}{@{}L{2.1cm}|L{6.1cm}|L{6.1cm}@{}}
\textbf{Method} & \textbf{How would it work here?} & \textbf{Does it fit? Why or why not?} \\
\hline
\textbf{Simple random sample} &
\cellans{2.0cm}{Number every home 1 to 9{,}000 and use a random number generator to pick 300.} &
\cellans{2.0cm}{Not practical: there is no list to number, and 300 homes scattered over 600
square miles means hundreds of hours of driving.} \\
\hline
\textbf{Stratified random sample} &
\cellans{2.0cm}{Split the homes into groups alike inside (say, old houses and new houses) and
take an SRS from each group.} &
\cellans{2.0cm}{No: nothing says which homes are old or new, so strata cannot be formed, and it still
means driving everywhere.} \\
\hline
\textbf{Cluster sample} &
\cellans{2.0cm}{Number the 150 map sections, randomly choose 5, and check every home in them
(5 $\times$ 60 = 300 homes).} &
\cellans{2.0cm}{Fits: each section is a mix, like a small copy of the county, and the inspector
drives to only 5 places.} \\
\end{tabular}
\end{scenariobox}

\vspace{0.08in}

\boxguard[12]
\begin{scenariobox}[3. Our choice]{navy}
\small
\begin{enumerate}[label=\textbf{\alph*.}, leftmargin=1.6em, itemsep=4pt]
  \item The best method for our card: \fillin{8.5cm}{a cluster sample (the map sections)}
  \item Defend it \emph{in context}: why does this method fit \emph{this} population better than
        the other two?
        \answerspace{1.8cm}{Each map section mixes old and new homes, owners and renters, so
        each one looks like the whole county. Checking every home in 5 sections gives about 300
        homes with only 5 trips. An SRS needs a list we do not have and would send the inspector
        all over 600 square miles.}
  \item Describe \emph{exactly} how to choose the sample, step by step.
        \answerspace{1.8cm}{Number the map sections 001 to 150. Use a random number generator to
        pick 5 different numbers from 1 to 150, skipping repeats. The inspector visits every
        home in those 5 sections and tests each smoke detector.}
\end{enumerate}
\end{scenariobox}

\vspace{0.08in}

\boxguard[8]
\begin{scenariobox}[4. Listen to the other tables]{navy}
\small
\renewcommand{\arraystretch}{1.2}
\begin{tabular}{@{}L{1.3cm}|L{4.0cm}|L{2.9cm}|L{5.8cm}@{}}
\textbf{Table} & \textbf{Their question} & \textbf{Best method} & \textbf{The reason, in one phrase} \\
\hline
\textbf{2} & \cellans{1.2cm}{StreamLine: mean rating of Daily Mix} & \cellans{1.2cm}{Simple random sample} &
\cellans{1.2cm}{Full list, free to reach anyone, nothing to stratify on.} \\
\hline
\textbf{3} & \cellans{1.2cm}{Northfield: mean commute distance} & \cellans{1.2cm}{Stratified, by work site} &
\cellans{1.2cm}{Commutes alike inside a site, very different between sites.} \\
\end{tabular}

\vspace{6pt}
\centering\bfseries
Strata are alike inside: sample from \fillin{1.4cm}{every} one.\par
Clusters each look like the whole: take \fillin{1.0cm}{all} of \fillin{1.3cm}{some} of them.
\end{scenariobox}

% ════════════════════════════════════════════════════════════════════════════
% TABLE 2 — STREAMLINE — SIMPLE RANDOM SAMPLE
% ════════════════════════════════════════════════════════════════════════════
\clearpage
\pageheader{Unit 3, Lesson 3.2}{Table Activity: Which Plan Fits? \quad --- \quad Table 2 \quad --- \quad Key}

\boxguard[10]
\begin{headlinebox}{sky}
\small
\textbf{Your table is the survey team.} Read your card, try all three sampling methods on it,
and choose the one that fits. Be ready to defend your choice to the class in one minute ---
\emph{in context}, using the words \emph{alike inside}, \emph{looks like the whole}, or
\emph{cost}.
\end{headlinebox}

\vspace{0.08in}

\boxguard[12]
\begin{scenariobox}[Table 2 card: Rating a New Feature on StreamLine]{navy}
\small
\textbf{Statistical question:} What is the mean rating, on a scale of 1 to 10, that StreamLine
subscribers give the app's new \emph{Daily Mix} feature?
\begin{itemize}[leftmargin=1.3em, itemsep=2pt, topsep=4pt]
  \item StreamLine, a music app, has 50{,}000 subscribers all over the country.
  \item \textbf{The company has a list of every subscriber}, numbered 00001 to 50000. The list
        records nothing else about them --- no age, no location, no listening habits.
  \item A survey can be sent inside the app to any subscriber \emph{at no cost}, and it reaches
        them wherever they live.
  \item The company wants \textbf{200 ratings}.
\end{itemize}
\end{scenariobox}

\vspace{0.08in}

\boxguard[10]
\begin{scenariobox}[1. Read the population]{navy}
\small
\begin{enumerate}[label=\textbf{\alph*.}, leftmargin=1.6em, itemsep=4pt]
  \item What are the individuals? What is the variable?
        \answerspace{1.1cm}{Individuals: StreamLine's 50{,}000 subscribers. Variable: the
        subscriber's rating of Daily Mix, 1 to 10 (quantitative).}
  \item Is there a list of every individual? What does it cost to reach one?
        \answerspace{1.3cm}{Yes --- a numbered list, 00001 to 50000. Reaching any subscriber
        costs nothing, and it costs the same for everyone, near or far.}
  \item Are there groups that are \emph{alike inside} on the response? Groups that each
        \emph{look like the whole} population?
        \answerspace{1.3cm}{None that we can use: the list records nothing that would split
        subscribers into groups alike inside, and there are no natural groups to use as
        clusters.}
\end{enumerate}
\end{scenariobox}

\clearpage

\boxguard[16]
\begin{scenariobox}[2. Try all three methods on your card]{navy}
\small
\renewcommand{\arraystretch}{1.2}
\begin{tabular}{@{}L{2.1cm}|L{6.1cm}|L{6.1cm}@{}}
\textbf{Method} & \textbf{How would it work here?} & \textbf{Does it fit? Why or why not?} \\
\hline
\textbf{Simple random sample} &
\cellans{2.0cm}{Use a random number generator to pick 200 different account numbers from 00001
to 50000 and send each the survey.} &
\cellans{2.0cm}{Fits: the list exists, every subscriber is reached for free, and nothing is known
that would let us do better.} \\
\hline
\textbf{Stratified random sample} &
\cellans{2.0cm}{Split subscribers into groups alike in how they would rate it (say, heavy and
light listeners) and take an SRS from each.} &
\cellans{2.0cm}{Cannot be done: the list records no age, location, or habits, so there is no way
to put subscribers into strata.} \\
\hline
\textbf{Cluster sample} &
\cellans{2.0cm}{Split the list into blocks of account numbers, randomly pick a few blocks, and
survey everyone in them.} &
\cellans{2.0cm}{No advantage: reaching anyone is already free, so clusters save nothing --- and
blocks of account numbers are not real groups.} \\
\end{tabular}
\end{scenariobox}

\vspace{0.08in}

\boxguard[12]
\begin{scenariobox}[3. Our choice]{navy}
\small
\begin{enumerate}[label=\textbf{\alph*.}, leftmargin=1.6em, itemsep=4pt]
  \item The best method for our card: \fillin{8.5cm}{a simple random sample (SRS)}
  \item Defend it \emph{in context}: why does this method fit \emph{this} population better than
        the other two?
        \answerspace{1.8cm}{StreamLine has a numbered list of all 50{,}000 subscribers and can
        survey any of them for free, so a cluster sample saves nothing. With no information on
        age or habits, there are no strata to form. An SRS gives every group of 200 subscribers
        the same chance to be the sample.}
  \item Describe \emph{exactly} how to choose the sample, step by step.
        \answerspace{1.8cm}{Use a random number generator to choose integers from 1 to 50{,}000,
        skipping repeats, until there are 200 different account numbers. Send the in-app survey
        to those 200 subscribers and average their ratings.}
\end{enumerate}
\end{scenariobox}

\vspace{0.08in}

\boxguard[8]
\begin{scenariobox}[4. Listen to the other tables]{navy}
\small
\renewcommand{\arraystretch}{1.2}
\begin{tabular}{@{}L{1.3cm}|L{4.0cm}|L{2.9cm}|L{5.8cm}@{}}
\textbf{Table} & \textbf{Their question} & \textbf{Best method} & \textbf{The reason, in one phrase} \\
\hline
\textbf{1} & \cellans{1.2cm}{Pine Ridge: homes with working smoke detectors} & \cellans{1.2cm}{Cluster, by map section} &
\cellans{1.2cm}{Each section is a mix like the county; only 5 trips.} \\
\hline
\textbf{3} & \cellans{1.2cm}{Northfield: mean commute distance} & \cellans{1.2cm}{Stratified, by work site} &
\cellans{1.2cm}{Commutes alike inside a site, very different between sites.} \\
\end{tabular}

\vspace{6pt}
\centering\bfseries
Strata are alike inside: sample from \fillin{1.4cm}{every} one.\par
Clusters each look like the whole: take \fillin{1.0cm}{all} of \fillin{1.3cm}{some} of them.
\end{scenariobox}

% ════════════════════════════════════════════════════════════════════════════
% TABLE 3 — NORTHFIELD FOODS — STRATIFIED (THE CRUX: a site is not a cluster)
% ════════════════════════════════════════════════════════════════════════════
\clearpage
\pageheader{Unit 3, Lesson 3.2}{Table Activity: Which Plan Fits? \quad --- \quad Table 3 \quad --- \quad Key}

\boxguard[10]
\begin{headlinebox}{sky}
\small
\textbf{Your table is the survey team.} Read your card, try all three sampling methods on it,
and choose the one that fits. Be ready to defend your choice to the class in one minute ---
\emph{in context}, using the words \emph{alike inside}, \emph{looks like the whole}, or
\emph{cost}.
\end{headlinebox}

\vspace{0.08in}

\boxguard[12]
\begin{scenariobox}[Table 3 card: Commuting to Northfield Foods]{navy}
\small
\textbf{Statistical question:} What is the mean one-way commute distance, in miles, for
Northfield Foods employees?
\par\vspace{4pt}
{\centering\footnotesize
\renewcommand{\arraystretch}{1.2}
\begin{tabular}{@{}l c l@{}}
\rowcolor{sky}
\textbf{Work site} & \textbf{Employees} & \textbf{Where most of them live} \\
\hline
Downtown plant & 1{,}200 & in the city, close to the plant \\
Riverbend farm, 40 miles out of the city & 400 & in the city --- they drive out \\
Lakeside office, in the suburbs & 200 & in the nearby suburbs \\
\end{tabular}\par}
\begin{itemize}[leftmargin=1.3em, itemsep=2pt, topsep=4pt]
  \item \textbf{Human Resources has a list of all 1{,}800 employees} and the site where each
        one works.
  \item Any employee can be surveyed by company email \emph{at no cost}.
  \item The company wants \textbf{90 responses}.
\end{itemize}
\end{scenariobox}

\vspace{0.08in}

\boxguard[10]
\begin{scenariobox}[1. Read the population]{navy}
\small
\begin{enumerate}[label=\textbf{\alph*.}, leftmargin=1.6em, itemsep=4pt]
  \item What are the individuals? What is the variable?
        \answerspace{1.1cm}{Individuals: Northfield Foods' 1{,}800 employees. Variable: one-way
        commute distance in miles (quantitative).}
  \item Is there a list of every individual? What does it cost to reach one?
        \answerspace{1.3cm}{Yes --- HR's list, with each employee's site. Email reaches anyone
        for free.}
  \item Are there groups that are \emph{alike inside} on the response? Groups that each
        \emph{look like the whole} population?
        \answerspace{1.3cm}{The sites are alike inside: Riverbend commutes are all long,
        downtown commutes mostly short. No site looks like the whole company.}
\end{enumerate}
\end{scenariobox}

\clearpage

\boxguard[16]
\begin{scenariobox}[2. Try all three methods on your card]{navy}
\small
\renewcommand{\arraystretch}{1.2}
\begin{tabular}{@{}L{2.1cm}|L{6.1cm}|L{6.1cm}@{}}
\textbf{Method} & \textbf{How would it work here?} & \textbf{Does it fit? Why or why not?} \\
\hline
\textbf{Simple random sample} &
\cellans{2.0cm}{Number all 1{,}800 employees, use a random number generator to pick 90, and
email them.} &
\cellans{2.0cm}{Works and is free, but by chance it could get too few Riverbend workers, and the mean
comes out too low.} \\
\hline
\textbf{Stratified random sample} &
\cellans{2.0cm}{Take an SRS from each site, in proportion: 60 downtown, 20 at Riverbend, 10 at
Lakeside.} &
\cellans{2.0cm}{Best: commutes are alike inside a site and differ between sites, so every site is in
the sample and the estimate varies less.} \\
\hline
\textbf{Cluster sample} &
\cellans{2.0cm}{Randomly pick one site and survey every employee there.} &
\cellans{2.0cm}{Wrong groups: one site is not a copy of the company (all Riverbend: mean far too
high), and email is free, so nothing is saved.} \\
\end{tabular}
\end{scenariobox}

\vspace{0.08in}

\boxguard[12]
\begin{scenariobox}[3. Our choice]{navy}
\small
\begin{enumerate}[label=\textbf{\alph*.}, leftmargin=1.6em, itemsep=4pt]
  \item The best method for our card: \fillin{8.5cm}{a stratified random sample, by work site}
  \item Defend it \emph{in context}: why does this method fit \emph{this} population better than
        the other two?
        \answerspace{1.8cm}{Commute distance depends on the site: Riverbend workers drive about 40
        miles, and downtown workers live nearby. Employees are alike inside a site, so sampling
        from every site makes sure each one is represented and gives a less variable estimate
        than an SRS of 90. A single site is not a copy of the company, so it is not a cluster.}
  \item Describe \emph{exactly} how to choose the sample, step by step.
        \answerspace{1.8cm}{From HR's list, number the downtown employees 1--1{,}200, Riverbend
        1--400, and Lakeside 1--200. Use a random number generator to pick 60, 20, and 10
        different numbers from those lists, skipping repeats. Email those 90 employees and
        combine all 90 commute distances.}
\end{enumerate}
\end{scenariobox}

\vspace{0.08in}

\boxguard[8]
\begin{scenariobox}[4. Listen to the other tables]{navy}
\small
\renewcommand{\arraystretch}{1.2}
\begin{tabular}{@{}L{1.3cm}|L{4.0cm}|L{2.9cm}|L{5.8cm}@{}}
\textbf{Table} & \textbf{Their question} & \textbf{Best method} & \textbf{The reason, in one phrase} \\
\hline
\textbf{1} & \cellans{1.2cm}{Pine Ridge: homes with working smoke detectors} & \cellans{1.2cm}{Cluster, by map section} &
\cellans{1.2cm}{Each section is a mix like the county; only 5 trips.} \\
\hline
\textbf{2} & \cellans{1.2cm}{StreamLine: mean rating of Daily Mix} & \cellans{1.2cm}{Simple random sample} &
\cellans{1.2cm}{Full list, free to reach anyone, nothing to stratify on.} \\
\end{tabular}

\vspace{6pt}
\centering\bfseries
Strata are alike inside: sample from \fillin{1.4cm}{every} one.\par
Clusters each look like the whole: take \fillin{1.0cm}{all} of \fillin{1.3cm}{some} of them.
\end{scenariobox}

\end{document}
